\chapter{\MakeUppercase{Introduction}}
\justifying
\section{Background}
This section explains the context of the project and sets the foundation. The content include:

\begin{itemize}
	\item Overview of the Problem Domain.
	\item Current Systems or Existing Solutions.
	\item Limitations of Existing Systems.
	\item Relevant Technologies and Concepts
	\item Theoretical Foundations
	\item Need for the Proposed System
\end{itemize}
\paragraph{Use of Acronym in the text:}The acronym Artificial Intelligence can be included as short \ac{ai} using latex command \verb|\ac{ai}|  in the text. Similarly the URL \ac{url} can be used as \verb|\ac{url}| in the text. Machine learning \ac{ml} can be used as \verb|\ac{ml}| in the text. Use of other acronyms \ac{2g}, \ac{3g}, \ac{csma} are  defined as \verb|\ac{2g}, \ac{3g}, \ac{csma}|, respectively.
\subsection{Sub-section heading 1}
\paragraph{} subsection content to be included 

\subsection{Sub-section heading 2}
\paragraph{} subsection content to be included 

\section{Motivation}
This section explains why the project was undertaken. The content include:
\begin{itemize}
	\item Problem Significance
	\item Practical Motivation
	\item Academic Motivation
	\item Technical Interest
	\item Innovation or Improvement Aspect
\end{itemize}

\section{Scope of the project}
The scope of the project section defines the boundaries and limitations. The content include:
\begin{itemize}
	\item Functional Scope
	\item Non-Functional Scope
	\item Technical Scope
	\item Project Boundaries 
		\begin{itemize}
			\item include assumption and constraints.
		\end{itemize}
\end{itemize}

